\documentclass{notes}

\hypersetup{
  pdfauthor={Your Name},
  pdftitle={Notes Template}
}

\title{信息论笔记与开卷资料模板}
\author{南极滑稽}
\date{\today}

\begin{document}

\begin{titlepage}
  \centering
  \vspace*{2.2cm}

  {\color{MidnightBlue}\rule{0.78\textwidth}{1.2pt}}\\[1.2cm]

  {\fontsize{30}{36}\bfseries\selectfont 信息论笔记与开卷资料}\\[0.5cm]
  {\fontsize{18}{22}\selectfont Information Theory Notes}\\[1.8cm]

  {\Large\textbf{南极滑稽}}\\[0.5cm]
  {\large \today}\\[1.6cm]

  \begin{minipage}{0.78\textwidth}
    \centering
    \setlength{\fboxsep}{8pt}
    \colorbox{MidnightBlue!5}{%
      \parbox{0.88\linewidth}{
        \centering
        信息是对不确定性的消除。\\
        -- 克劳德·香农
      }
    }
  \end{minipage}\\[1.8cm]

  {\color{MidnightBlue}\rule{0.78\textwidth}{1.2pt}}

  \vfill
\end{titlepage}

\pagenumbering{arabic}
\tableofcontents
\newpage

\section{使用说明}

本模板适合整理课堂笔记、复习提纲和开卷考试资料。建议每节按照“定义 $\to$ 定理 $\to$ 公式 $\to$ 例题 $\to$ 易错点”的顺序整理。

\begin{keypoint}
开卷资料应优先记录可快速检索的信息：核心结论、适用条件、常见变形、典型题型和容易漏掉的限制条件。
\end{keypoint}

\begin{hintBox}
需要更紧凑的排版时，可以把某些总结页放入 \verb|multicols| 环境，例如 \verb|\begin{multicols}{2} ... \end{multicols}|。
\end{hintBox}

\section{基础概念}

\begin{definition}[熵]
离散随机变量 $X$ 的熵定义为
\[
  H(X) = -\sum_{x \in \mathcal{X}} p(x) \log p(x).
\]
若没有特别说明，信息论中常取 $\log = \log_2$，单位为 bit。
\end{definition}

\begin{formulaBox}
常用关系式：
\begin{align*}
  H(X,Y) &= H(X) + H(Y \mid X), \\
  I(X;Y) &= H(X) - H(X \mid Y), \\
  I(X;Y) &= H(X) + H(Y) - H(X,Y), \\
  \KL(P\|Q) &= \sum_x P(x) \log \frac{P(x)}{Q(x)}.
\end{align*}
\end{formulaBox}

\begin{theorem}[互信息非负]
对任意离散随机变量 $X,Y$，有
\[
  I(X;Y) \ge 0.
\]
等号成立当且仅当 $X$ 与 $Y$ 相互独立。
\end{theorem}

\begin{proof}
由 Gibbs 不等式，
\[
  I(X;Y) = \KL(P_{XY}\|P_XP_Y) \ge 0.
\]
当且仅当 $P_{XY}=P_XP_Y$ 时取等号，即 $X$ 与 $Y$ 独立。
\end{proof}

\begin{warningBox}
不要混淆 $H(X\mid Y)=0$ 与 $I(X;Y)=0$。前者表示 $Y$ 完全决定 $X$，后者表示 $X$ 与 $Y$ 没有统计相关性。
\end{warningBox}

\section{题型整理}

\begin{methodBox}[title={计算熵和互信息的步骤}]
\begin{enumerate}
  \item 写出联合分布表或概率质量函数。
  \item 求边缘分布 $P_X, P_Y$。
  \item 选择最短路径：用 $I(X;Y)=H(X)-H(X\mid Y)$ 或 $I(X;Y)=H(X)+H(Y)-H(X,Y)$。
  \item 检查单位和对数底数。
\end{enumerate}
\end{methodBox}

\begin{example}
若 $X\sim\mathrm{Bernoulli}(p)$，则
\[
  H(X)=h_2(p)=-p\log p-(1-p)\log(1-p).
\]
特别地，$p=1/2$ 时 $H(X)=1$ bit。
\end{example}

\section{开卷速查页示例}

\begin{multicols}{2}
\subsection*{核心公式}
\begin{itemize}
  \item 条件熵：$H(X\mid Y)=H(X,Y)-H(Y)$
  \item 链式法则：$H(X_1^n)=\sum_{i=1}^n H(X_i\mid X_1^{i-1})$
  \item 数据处理不等式：若 $X\to Y\to Z$，则 $I(X;Z)\le I(X;Y)$
\end{itemize}

\subsection*{常见判断}
\begin{itemize}
  \item 独立 $\Rightarrow I(X;Y)=0$
  \item 函数关系 $X=f(Y) \Rightarrow H(X\mid Y)=0$
  \item 条件会降低熵：$H(X\mid Y)\le H(X)$
\end{itemize}
\end{multicols}

\end{document}
